\chapter{Estudo analítico de cliques}
\label{ap:cliques}

Este estudo compara as ações necessárias para consultar alternativas de transporte público e selecionar uma rota no PaP, Google Maps, Moovit, Cittamobi e Quanto Tempo Falta. São considerados endereço digitado, favorito comum, atalhos Casa e Trabalho e local selecionado no mapa. As contagens foram obtidas pelo grupo em 25 e 26 de outubro de 2026; não integram o parecer de Thiago nem constituem avaliação de usabilidade com passageiros.

\papsection{Procedimento e condições}
O aparelho de referência foi um Samsung Galaxy S24 Ultra (SM-S928B), Android 16, com tela de 1080 por 2340 pixels. Foram utilizados Google Maps \nolinkurl{26.43.00.1002020710}, Moovit \nolinkurl{5.201.6.1811}, Cittamobi \nolinkurl{7.2.164}, Quanto Tempo Falta \nolinkurl{v2.0.2} e PaP no Chrome \nolinkurl{155.0.8059.84}. Os manifestos das sessões identificam a compilação do PaP e as capturas correspondentes.

A contagem começa na tela inicial, com localização disponível, permissões respondidas e favoritos cadastrados. Cada toque em um controle, toque longo, gesto de rolagem ou arrasto do mapa conta como uma ação de navegação, com o mesmo peso. A tabela soma toques e rolagens; os registros preservam a distinção entre esses gestos. A digitação do endereço constitui um episódio separado, sem contar cada tecla. Instalação, autenticação e cadastro dos favoritos ficam fora da tarefa. O marco final é a abertura da rota selecionada ou, na consulta de favoritos de parada, do trajeto de uma linha. Não se conta o início do acompanhamento durante o deslocamento.

Em São Paulo, a origem simulada foi a região do MASP, $(-23{,}561458; -46{,}656550)$, e o destino pesquisado foi Praça da Sé. O favorito de ponto usado nas medições do PaP foi \texttt{800016588}, na região da Sé/Praça Clóvis Beviláqua. No Quanto Tempo Falta, foi utilizada Santos, com origem $(-23{,}9675; -46{,}3336)$ e destino Praça Mauá. Os destinos resolvidos pelos aplicativos não são rigorosamente iguais; comparam-se ações de interface em áreas atendidas, não a qualidade dessas viagens. A localização foi simulada e o GPS restaurado ao concluir os testes.

\papsection{Resultados consolidados}
A Tabela~\ref{tab:cliques-celular} apresenta o total de ações até abrir a rota selecionada. Nos favoritos do Cittamobi e Quanto Tempo Falta, o marco é abrir o trajeto da linha: o percurso analítico informado pelo grupo é abrir Favoritos, escolher a parada e abrir a linha, totalizando três ações. Esses últimos valores representam a reconstrução do percurso descrito, sem nova medição integral no aparelho. Os resultados finais dessas tarefas são distintos e não sustentam uma classificação única dos aplicativos. As contagens dependem do estado inicial e não demonstram um mínimo universal.

\begin{table}[H]
\centering\small\SingleSpacing
\caption{Total de ações até abrir a rota ou o trajeto da linha.}
\label{tab:cliques-celular}
\setlength{\tabcolsep}{4pt}\renewcommand{\arraystretch}{1.18}
\begin{tabular}{>{\raggedright\arraybackslash}p{0.22\linewidth}>{\centering\arraybackslash}p{0.12\linewidth}>{\centering\arraybackslash}p{0.15\linewidth}>{\centering\arraybackslash}p{0.08\linewidth}>{\centering\arraybackslash}p{0.11\linewidth}>{\centering\arraybackslash}p{0.17\linewidth}}
\hline
Aplicativo & Endereço & Favorito comum & Casa & Trabalho & Local no mapa \\\hline
PaP & 4 & 3 & 2 & 2 & 3 \\[3pt]
Google Maps & 4 & 4 & 2 & 4 & 3 \\[3pt]
Moovit & 4 & 3 & 3 & 3 & 6 \\[3pt]
Cittamobi & 5 & 3 & 3 & 3 & 8 \\[3pt]
Quanto Tempo Falta & 4 & 3 & 3 & 3 & \textemdash \\\hline
\end{tabular}
\par\smallskip\begin{minipage}{0.96\linewidth}\footnotesize
Fonte: Elaboração do grupo a partir dos registros de execução no celular em 25 e 26 out. 2026 e da reconstrução analítica do percurso de favoritos do Cittamobi e Quanto Tempo Falta.
\end{minipage}
\end{table}

\textbf{Local no mapa.} No PaP, foram três ações: toque longo, confirmação do destino e seleção da rota a pé de 358 m. No Google Maps, foram três com transporte público mantido e quatro quando foi preciso mudar o modal. No Moovit, foram seis ações: cinco toques e um arrasto para posicionar o mapa; o fechamento de um anúncio elevou a execução observada a sete. A coluna mantém as seis ações de navegação, explicitando essa interrupção. Os destinos e enquadramentos foram diferentes; deslocamentos adicionais do mapa aumentam os totais.

No Cittamobi, foram oito ações desde o mapa inicial: abrir o menu, acessar ``Traçar Rotas'', escolher um destino do histórico para entrar no planejador, editar o destino, acessar ``Selecione no mapa'', arrastar o mapa, confirmar em ``OK'' e selecionar uma rota. Foram sete toques e um arrasto, com os detalhes da linha 147 abertos em São Bernardo do Campo. A contagem descreve esse percurso observado, com histórico disponível, sem demonstrar um mínimo universal. No Quanto Tempo Falta, o traço indica que a seleção de local livre não foi encontrada nas telas testadas em Santos; não significa zero ações nem ausência comprovada em todas as versões.

\textbf{PaP.} O favorito comum exigiu abrir Favoritos, usar o ponto como destino e selecionar a rota. Casa e Trabalho dispensaram a abertura da lista, reduzindo uma ação até a rota. As três medições utilizaram o mesmo ponto e não exigiram rolagem. A versão atual também permite salvar coordenadas de um local, inclusive pelo mapa ou pela localização atual, e associá-las a Casa ou Trabalho. As contagens da tabela foram medidas com favorito de ponto; o fluxo completo com favorito de coordenadas ainda precisa ser recontado no celular de referência. Por endereço, o percurso foi abrir a busca, digitar Praça da Sé, escolher a primeira sugestão, selecionar o ponto mostrado e selecionar a rota. Foram quatro ações até a rota da linha 715M-10 no mapa, sem rolagem e com um episódio de digitação. O destino selecionado foi o ponto \texttt{800016609}, Praça da Sé, 21.

\textbf{Google Maps.} O transporte público estava mantido como modal. Por endereço, o percurso foi abrir a busca, escolher a sugestão, pedir Rotas e selecionar a alternativa. Pelo favorito comum, foi necessário abrir Salvos, escolher o lugar diretamente visível, pedir Rotas e selecionar a alternativa. Casa estava acessível na tela inicial; Trabalho foi acessado por Salvos. Se o modal inicial for automóvel, a mudança para transporte público acrescenta um toque ao total. Os destinos pessoais de Casa e Trabalho foram preservados, sem divulgação dos endereços; essas contagens descrevem os acessos examinados, não qualquer disposição dos atalhos.

\textbf{Moovit.} Por endereço, foram necessários busca, sugestão, pedido de rotas e seleção da alternativa. O favorito comum e os atalhos Casa e Trabalho abriram diretamente as alternativas, mas uma rolagem foi necessária para alcançá-los na tela inicial. A seleção de rota acrescentou um toque. Os três favoritos apontavam para Praça da Sé. Conteúdos dinâmicos e anúncios podem acrescentar ações; a tabela apresenta os menores percursos concluídos, sem essas interrupções.

\textbf{Cittamobi e Quanto Tempo Falta.} O total de três ações considera o acesso pela lista de favoritos, com a parada visível após abrir essa lista. Casa e Trabalho são nomes atribuídos a favoritos comuns, seguindo a mesma sequência; não representam atalhos especiais fixos. Abrir a linha mostra seu trajeto, sem comprovar planejamento da localização atual até o ponto salvo. Se o favorito estiver fora da área visível, as rolagens necessárias acrescentam ações ao total. O acesso direto a um cartão já exposto na tela inicial constitui outro estado inicial e não substitui essa sequência.

\papsection{Interpretação e alcance}
Nas condições observadas, Casa e Trabalho reduziram uma ação em relação ao favorito comum do PaP. Google Maps e Moovit também apresentaram atalhos de acesso direto, condicionados à posição dos controles. Ao incluir a rolagem, o favorito comum do Moovit totalizou as mesmas três ações observadas no PaP. Os totais dos outros dois aplicativos descrevem consulta de linhas e não são equivalentes ao planejamento de viagem.

As contagens por favorito dependem de sua posição na interface. Listas maiores podem exigir rolagens, abertura de listas ou busca; cada toque e cada gesto de rolagem acrescenta uma ação. O ensaio com cartões fora da tela não constitui teste de listas com 5, 10 ou 20 favoritos. Não foi medido o efeito dessas grandes quantidades, e os valores não devem ser generalizados para esse cenário.

O estudo mede ações de navegação, não tempo de resposta, esforço cognitivo, acessibilidade, precisão operacional ou qualidade das rotas. As diferenças de destinos, objetos salvos e estados iniciais limitam a comparação. A avaliação com passageiros continua necessária para verificar compreensão, eficácia e satisfação.

O modelo KLM de \citet{card1980} relaciona operações de interação ao tempo de execução de tarefas, incluindo preparação mental e resposta do sistema. A contagem aqui adotada não aplica integralmente esse modelo: atribui o mesmo peso a toques e rolagens e não estima seus tempos. \citet{laubheimer2019} também ressalta que o número de cliques isolado não determina a facilidade de uso. Assim, menos ações constituem um indicador limitado de economia de navegação, não uma comprovação da qualidade global do produto.

Os registros que sustentam a tabela, com versões, ações e hashes das capturas, estão em \mbox{\nolinkurl{estudo_cliques_evidencias/celular/}}; as sessões de Casa e Trabalho estão no subdiretório \mbox{\nolinkurl{atalhos-casa-trabalho/}}. Essas evidências permitem conferir as contagens e as condições de cada tarefa.

O complemento de 26 de outubro está no subdiretório \mbox{\nolinkurl{favoritos-2026-10-26/}}.

\section{Protocolo ampliado de avaliação da eficiência e compreensão}
\label{sec:ihc-ampliado}
A ampliação proposta pergunta se o acesso contextualizado pelo QR Code facilita a consulta e a escolha do deslocamento sem comprometer a compreensão do passageiro. Trata-se de protocolo futuro, distinto das contagens executadas pelo grupo e da avaliação exploratória por usuários convidados descrita nos resultados. A relação entre eficácia, eficiência e satisfação orienta a avaliação contextualizada \citep{iso924111}; o KLM de \citet{card1980} não é aplicado integralmente às contagens existentes. A recomendação de pequenas rodadas qualitativas \citep{nielsen2000} orienta a identificação de problemas, sem conferir representatividade estatística à amostra.

A ISO 25062:2025 organiza o conteúdo de relatórios conforme a abordagem de avaliação adotada \citep{iso25062}. Neste protocolo, essa orientação fundamenta a identificação explícita de participantes, tarefas, contexto, procedimento e limites das comparações, sem declarar conformidade integral à norma. A distinção entre sucesso, tempo, erros, eficiência e aprendizagem discutida por Tullis e Albert~(\citeyear{tullis2013}) fundamenta o uso combinado dessas medidas: a contagem de ações não substitui a conclusão correta nem a compreensão do percurso.

A auditoria preservou as Seções~\ref{sec:cliques-condicoes} a~\ref{sec:cliques-interpretacao} e a Tabela~\ref{tab:cliques-celular}. Nos 36 manifestos inventariados, 171 capturas conferiram com seus hashes e três foram omitidas por privacidade. Das referências dos consolidados, 24 manifestos conferiram; uma tentativa excluída do Moovit não está no acervo exportado. Favoritos do Cittamobi e Quanto Tempo Falta continuam como reconstruções de consulta de linha, distintas do planejamento por destino. Cidades, destinos, versões, anúncios e estados iniciais limitam comparações. No mapa, o PaP exibiu caminhada e outras aplicações exibiram ônibus. A captura final do Quanto Tempo Falta por endereço ainda mostra a lista após o toque de seleção: o marco de detalhe requer revisão, mantendo-se a contagem original. A matriz e as ressalvas completas acompanham os instrumentos de apoio.

As tarefas são classificadas em: T1, consultar próximos ônibus; T2, planejar por destino; T3, reutilizar destino favorito; T4, corrigir seleção; T5, planejar com baldeação; e T6, identificar origem por QR Code. Há percursos observados de T1, T2 e T3, com distinção entre favorito de parada para consultar linhas e favorito usado como destino. Correções ocasionais não constituem ensaio controlado de T4. Nem T5 nem a comparação controlada de T6 com a identificação manual da origem estão concluídas neste estudo; demonstrações e navegações exploratórias não substituem essas medições. A matriz detalhada, os instrumentos e a auditoria estão em \nolinkurl{estudo_ihc_ampliado/}.

\textbf{Delineamento proposto.} Cada participante realizará duas condições no PaP: MAN, identificação manual da origem; QR, identificação por leitura de QR Code. A hipótese de eficiência é que QR exija menos ações e menor duração; a hipótese de compreensão é que a identificação contextualizada não prejudique o reconhecimento de linha, sentido, embarque e baldeação. Ambas permanecem sem resultado. A variável independente é o método de identificação; as dependentes são ações, tempos, conclusão correta, erros, recuperação, ajuda, compreensão e facilidade percebida. Serão alternadas as ordens MAN$\rightarrow$QR e QR$\rightarrow$MAN entre participantes, com diferença máxima de um participante por ordem, registrando possível aprendizagem na repetição.

Serão mantidos parada e destino públicos, versão do PaP, dispositivo, navegador, permissões, rede, horário de partida e conjunto de dados. Antes da coleta, serão conferidos os itinerários admissíveis e as restrições do cenário. O gabarito aceitará múltiplas soluções viáveis: não exigirá a mesma linha, a primeira sugestão ou o menor tempo, salvo quando isso fizer parte expressa da tarefa. Para cada alternativa serão verificadas a continuidade entre trechos, a compatibilidade temporal, o sentido, as caminhadas e as baldeações; uma solução não antecipada será avaliada pelos mesmos critérios, com justificativa registrada. O cenário principal usará dados congelados e identificados, declarados como controlados, com restauração do estado e tratamento idêntico de cache entre condições; uma eventual avaliação com dados operacionais será separada. O QR Code deverá abrir a mesma aplicação e codificar a parada correta, com servidor acessível pelo celular. Não é necessário domínio permanente. Na condição MAN, não haverá origem pré-preenchida; na condição QR, não haverá leitura antecipada. A equivalência será registrada antes da coleta: mesmos objetivos e critérios de sucesso, limites de caminhada, quantidade de baldeações, janela temporal, oferta de alternativas válidas, estado inicial e familiaridade com o trajeto. Cenários que não atendam a esses critérios serão analisados separadamente, sem comparação direta de suas contagens ou tempos.

\textbf{Marcos e métricas.} Após ler a tarefa, o participante partirá do mesmo estado neutro do dispositivo, sem aplicação ou câmera abertas. O sinal do moderador define $t_0$; o acesso MAN inclui abrir o PaP e indicar a parada, e o acesso QR inclui abrir o leitor, enquadrar, decodificar e abrir o endereço. Registrar-se-ão origem confirmada ($t_o$), destino confirmado ($t_d$), primeira alternativa visível que atenda ao cenário ($t_a$), término do cálculo sinalizado ($t_c$), seleção da alternativa ($t_s$) e abertura estável de seus detalhes ($t_f$). Os indicadores de duração são diferenças em relação a $t_0$; $t_s-t_a$ caracteriza a escolha após a oferta de alternativa. Se não houver sinal confiável de término do cálculo, $t_c$ ficará ausente, sem substituí-lo pela primeira resposta HTTP. Falhas e interrupções serão registradas com sua causa, sem imputar duração de sucesso.

Cada tarefa registrará separadamente toques, toques prolongados, rolagens, arrastos e episódios de digitação. Um gesto contínuo conta uma unidade; eventos intermediários de rolagem não contam como gestos adicionais. Digitação compreende entrada contínua no mesmo campo, encerrada por submissão ou mudança de foco. A soma de ações manterá a convenção da Seção~\ref{sec:convencao-acoes}, com digitação em coluna própria. Conclusão correta exige origem e destino do cenário e alternativa compatível com suas restrições; erros são seleções incompatíveis com a tarefa, não simples exploração. Recuperação compreende ações para corrigir erro identificado; pedidos de ajuda e orientações fornecidas serão separados. Ao final, o participante explicará linha, sentido, embarque e local de baldeação, quando aplicável; cada item será classificado como correto, parcial, incorreto ou não aplicável, segundo gabarito prévio. A facilidade percebida será perguntada após cada tentativa pelo SEQ, em sete pontos, de muito difícil a muito fácil \citep{sauro2012seq}; sua formulação em português será tratada como adaptação de trabalho, sem alegar validação psicométrica.

\textbf{Participantes e execução.} Propõe-se rodada qualitativa com cinco a oito participantes externos ao desenvolvimento, registrando separadamente familiaridade com transporte e com aplicativos de navegação. O roteiro terá acolhimento e consentimento, prática neutra sem revelar o caminho avaliado, execução na ordem MAN$\rightarrow$QR ou QR$\rightarrow$MAN, explicação do percurso e SEQ após cada condição. A prática terá duração e instruções iguais; não haverá demonstração da solução nem correção orientada entre MAN e QR. Serão registrados ordem, tentativa e experiência prévia. Como repetir o mesmo cenário pode favorecer a segunda condição, os resultados serão examinados também por ordem MAN$\rightarrow$QR e QR$\rightarrow$MAN; o contrabalanceamento reduz esse viés, mas não comprova sua eliminação. Se um estudo posterior usar trajetos diferentes para reduzir memorização, deverá conferir previamente sua equivalência e contrabalancear também a associação entre trajeto e condição; seus dados não serão misturados aos pares do mesmo cenário. Tarefas complementares curtas avaliarão consulta de próximos ônibus, favorito público previamente cadastrado, correção de um destino incorreto preparado pelo moderador e viagem com uma baldeação. Não serão usados endereços pessoais. A contagem experimental inclui publicidade e interrupções encontradas, discriminadas por tipo. A coleta depende de autorização própria e não foi iniciada nesta ampliação.

\textbf{Telemetria, análise e limites.} Os estados e traces existentes permitem localizar origem, destino, mudanças, retornos, cancelamentos e seleção, desde que a captura esteja completa e a semântica da versão seja conferida. Consentimento tardio, reinício da página e envio em lotes impedem tratar o primeiro evento recebido como início da tarefa; o relógio monotônico reinicia com a página. Resposta de API não comprova alternativa utilizável, renderização ou conclusão correta. Ajuda, intenção, compreensão e SEQ exigem observação ou pergunta própria. Serão separados comportamento observado, interpretação possível e dificuldade confirmada pelo relato ou por erro verificado em tarefa definida.

Para futuras repetições comparáveis, serão apresentados número de tentativas válidas e incompletas, valores individuais, mediana, mínimo, máximo e intervalo interquartil, com método de quartis declarado. Diferenças QR menos MAN serão calculadas por participante apenas para pares válidos do mesmo cenário, mantendo falhas visíveis na análise de eficácia. Uma observação isolada será apresentada como tal, sem inferência populacional; não serão combinadas reconstruções analíticas, versões ou tarefas distintas. Não se propõem testes de significância nem percentuais de melhoria com os dados atuais. Permanecem pendentes o experimento nas ordens MAN$\rightarrow$QR e QR$\rightarrow$MAN, medidas de tempo e compreensão, SEQ, recuperação controlada, baldeação, favorito de coordenadas e listas maiores. Menos ações exigem conclusão correta e compreensão preservada.

